\section{Conclusion}

We have argued that hypervisor--OS coupling is primarily a software-architecture
issue and that it can be addressed by making the dependency explicit through a
substrate interface contract. By restructuring \axvisor{} around such a
contract, the same virtualization core runs on \arceos{}, \asterinas{}, and
Linux, while a KVM-compatible ABI keeps the guest side decoupled from the
substrate choice. Booting a bare-metal Linux guest on both the \asterinas{} and
Linux substrates demonstrates that the contract is sufficient for real
workloads, and the measured adaptation effort and overhead indicate that the
cost of portability is modest. The broader lesson is that hypervisors need not
be bound to a single host OS; with the right service boundary, they can be
engineered as portable system components and retargeted across kernels with
limited effort.
